% Compacte DV2-resultatentekst, opgesteld op 2026-06-05.
% Bron: bevroren DV2-bundels 01b en 02b, include_english=False,
% USE_FROZEN_ADVIES=True, USE_171_SCOPE_KABINET=True.

\subsection{Betrouwbaarheid van de meting (DV2)}
\label{subsec:dv2-betrouwbaarheid-compact}

\noindent DV2 toetst niet of de AI-pijplijn foutloos is, maar of zij betrouwbaar
genoeg is om tekstuele relaties op corpusschaal te meten. De leidende peildatum is
5 juni 2026. De validatie gebruikt alleen Nederlandstalige documenten
(\texttt{include\_english=False}) en de bevroren 171-Kaderwet-scope. Daarmee rust
de adviesvalidatie op 1.374 adviesrapporten met 22.286 aanbevelingen en 17.927
probleemdefinities. De kabinetsreactievalidatie rust op 419 in-scope
kabinetsreactiedocumenten met samen 13.165 advieselementen. Deze noemers meten
niet hetzelfde als de DV1-contextcijfers: DV2 gaat over meetkwaliteit van de
AI-output, niet over de volledige corpusopbouw.

\noindent Het gebruiksoordeel is daarom belangrijker dan een enkel eindcijfer.
Een hoge score betekent hier niet dat elke individuele passage foutloos is
gecodeerd, maar dat de betreffende variabele sterk genoeg is om op geaggregeerd
niveau mee te tellen. Omgekeerd betekent een lage score niet dat de hele pijplijn
onbruikbaar is. Zij wijst vooral aan welke subvelden buiten de harde analyse
moeten blijven. DV2 fungeert daarmee als kwaliteitsfilter voor DV3: alleen
onderdelen die groen of voorzichtig oranje scoren dragen de inhoudelijke
resultaten.

\noindent De adviesextractie is het sterkste onderdeel van de pijplijn. Op de
golden set van 43 gekoppelde en gevalideerde documenten worden aanbevelingen
inhoudelijk goed teruggevonden (F1 = 0,910). Voor probleemdefinities is de
semantische score nog iets hoger (F1 = 0,915; 95\%-interval ongeveer
0,912--0,919). De lagere item-F1-scores laten zien dat de exacte afbakening van
losse items moeilijker is dan het terugvinden van de inhoudelijke strekking. Dat
is voor deze thesis aanvaardbaar, omdat DV3 patronen op corpusschaal rapporteert.
Twee subvelden zijn daarvoor niet sterk genoeg: consultatie-actoren (F1 = 0,216)
en beleidslogica-koppelingen (strikte F1 = 0,012). Die worden daarom niet als
harde kwantitatieve variabelen gebruikt.

\begin{table}[htbp]
\centering
\caption{Kernscores van de DV2-validatie, peildatum 5 juni 2026.}
\label{tab:dv2-compact-kernscores}
\small
\begin{tabularx}{\textwidth}{@{}>{\raggedright\arraybackslash}p{3.8cm}
>{\raggedright\arraybackslash}p{2.8cm}
>{\raggedright\arraybackslash}p{2.8cm}
>{\raggedright\arraybackslash}X@{}}
\toprule
Onderdeel & Noemer & Score & Interpretatie \\
\midrule
Aanbevelingen, semantische dekking & 43 documenten & F1 = 0,910 & Hard bruikbaar op corpusschaal \\
Aanbevelingen, itemniveau & 43 documenten & F1 = 0,828 & Bruikbaar, maar minder precies per item \\
Probleemdefinities, semantische dekking & 43 documenten & F1 = 0,915 & Hard bruikbaar op corpusschaal \\
Probleemdefinities, itemniveau & 43 documenten & F1 = 0,768 & Bruikbaar met nuance \\
Consultatie-actoren & 43 documenten & F1 = 0,216 & Niet hard gebruiken \\
Beleidslogica-koppelingen & 43 documenten & F1 = 0,012 & Niet hard gebruiken \\
Documentclassificatie & 101 documenten & 80/101 = 79,2\% & Geschikt voor corpusordening \\
Metadata canonical & 101 documenten & 61/101 = 60,4\% & Niet als harde individuele bron gebruiken \\
\bottomrule
\end{tabularx}
\end{table}

\noindent De kabinetsreactie-agent is bruikbaar, maar duidelijk zwakker dan de
adviesextractie. In de in-scope golden/silver-set vallen 87 van de 88 elementen
binnen de actuele afbakening. De exacte match op grove verwerkingslabels is
46/87, oftewel 52,9\%. Dat percentage moet niet als corpusbrede accuratesse
worden gelezen, omdat de testset per voorspeld label is gebalanceerd. De agent
herkent wel meestal dat een element behandeld wordt: 68 van de 77 behandelde
elementen worden als behandeld herkend. Niet-behandeling is minder zeker: 5 van
de 10 niet-behandelde elementen worden goed herkend. De belangrijkste inhoudelijke
verwarring ligt tussen substantiele en retorische verwerking. Juist daarom zijn
deze labels geschikt voor geaggregeerde patronen, maar niet voor harde uitspraken
over losse advieselementen zonder broncontrole.

\noindent Ook de aard van de referentie telt mee. De kabinetsreactievalidatie is
een golden/silver-set: zij is zorgvuldig gecontroleerd, maar niet hetzelfde als een
volledige menselijke dubbelcodering door meerdere onafhankelijke codeurs. Daardoor
is de score geschikt om de zwakke en sterke plekken van de agent te beoordelen,
maar niet om volledige reproduceerbaarheid te claimen. Voor de interpretatie van
DV3 betekent dit dat de brede verdeling van labels informatief is, terwijl
voorbeelden en grensgevallen terug naar de bron moeten.

\begin{table}[htbp]
\centering
\caption{Stabiliteit en dekking van de DV2-validatie.}
\label{tab:dv2-compact-stabiliteit}
\small
\begin{tabularx}{\textwidth}{@{}>{\raggedright\arraybackslash}p{4.0cm}
>{\raggedright\arraybackslash}p{3.3cm}
>{\raggedright\arraybackslash}X@{}}
\toprule
Controle & Uitkomst & Betekenis \\
\midrule
Adviesextractie hertest & 10 documenten x 3 runs; PEO 0,917--1,000 & Rapportlabels zijn grotendeels stabiel; itemsets blijven minder stabiel \\
Classificatie/metadata hertest & 100 documenten x 5 runs; alpha kritiek 0,935; metadata 0,959 & Sterk genoeg voor corpusordening en controle \\
Kabinetsreactie hertest & 10 cases x 3 runs; agreement 0,2537; alpha 0,1847 & Losse verwerkingslabels zijn niet stabiel genoeg als harde claim \\
Kabinetsreactie-golden-set & 46/87 exact; 68/77 behandeld; 5/10 niet-behandeld & Breed patroon bruikbaar, precieze labels voorzichtig lezen \\
Traceerbaarheid & 10.580/13.165 elementen met minstens drie trace-warnings & Broncontrole blijft nodig; warnings zijn niet automatisch inhoudelijke fouten \\
\bottomrule
\end{tabularx}
\end{table}

\noindent De stabiliteitscontrole bevestigt dit beeld. Classificatie en metadata
zijn reproduceerbaar genoeg voor ordening van het corpus. Adviesextractie blijft
stabiel genoeg voor de hoofdvariabelen, al is de exacte itemset bij herhaling
minder vast. De kabinetsreactie-agent is op elementniveau het minst stabiel:
herhaalde runs leveren vaak andere fijne labels op. Dat maakt de meting niet
onbruikbaar, maar wel beperkt. De juiste toepassing is aggregatie: aandelen en
patronen op corpusniveau, gelezen met de onzekerheid uit deze validatie.

\noindent Traceerbaarheid verandert dit oordeel niet, maar bepaalt wel hoe de
uitkomsten gelezen moeten worden. Veel kabinetsreactie-elementen hebben technische
trace-warnings. Dat zijn waarschuwingen over bronondersteuning, geen automatisch
bewijs dat het inhoudelijke label fout is. Zij maken wel duidelijk dat losse
casusclaims nooit alleen op het tabelresultaat mogen rusten. De AI-output is dus
een schaalinstrument: zij maakt het corpus systematisch meetbaar, waarna
broncontrole nodig blijft voor gevoelige interpretaties.

\begin{table}[htbp]
\centering
\caption{Verantwoord gebruik van de AI-output in de thesis.}
\label{tab:dv2-compact-gebruik}
\small
\begin{tabularx}{\textwidth}{@{}>{\raggedright\arraybackslash}p{3.0cm}
>{\raggedright\arraybackslash}p{2.2cm}
>{\raggedright\arraybackslash}X
>{\raggedright\arraybackslash}p{3.6cm}@{}}
\toprule
Onderdeel & Oordeel & Wel gebruiken voor & Niet gebruiken voor \\
\midrule
Aanbevelingen en probleemdefinities & Groen & Corpusbrede analyse van advieselementen & Ongecontroleerde claims over elke losse formulering \\
Documentclassificatie & Groen/oranje & Ordening en controle van het corpus & Definitieve juridische of inhoudelijke duiding per document \\
Metadata & Oranje & Context, ordening en broncontrole & Harde individuele uitspraken zonder controle \\
Kabinetsreactielabels & Oranje & Geaggregeerde patronen in zichtbare verwerking & Losse elementclaims zonder broncontrole \\
Consultatie-actoren en beleidslogica & Rood & Hoogstens verkennende context & Harde kwantitatieve conclusies \\
\bottomrule
\end{tabularx}
\end{table}

\noindent Het antwoord op DV2 is daarom positief, maar begrensd. De AI-pijplijn is
betrouwbaar genoeg om op corpusschaal aanbevelingen, probleemdefinities,
documenttypen en brede patronen in kabinetsreacties te meten. Zij is niet
betrouwbaar genoeg om zonder broncontrole harde uitspraken te doen over
individuele kabinetsreactielabels, consultatie-actoren of
beleidslogica-koppelingen. Voor DV3 betekent dit dat de hoofdconclusies op
geaggregeerd niveau worden gelezen. Losse voorbeelden blijven alleen overtuigend
wanneer zij teruggaan naar de bronpassage.
